\documentclass[11pt,a4paper]{ctexart}

\usepackage{amsmath,amssymb,bm}
\usepackage{geometry}
\usepackage{enumitem}
\usepackage{fancyhdr}
\usepackage{lastpage}
\usepackage{xcolor}
\usepackage{tcolorbox}
\usepackage{tikz}
\usepackage{titlesec}
\usepackage{microtype}
\usetikzlibrary{arrows.meta,calc}

\geometry{left=2.15cm,right=2.15cm,top=2.0cm,bottom=2.0cm}
\setlength{\parindent}{2em}
\setlength{\parskip}{0.25em}
\linespread{1.12}

\definecolor{mainblue}{RGB}{39,91,150}
\definecolor{lightblue}{RGB}{232,241,252}
\definecolor{darkgray}{RGB}{65,65,65}
\definecolor{lightgray}{RGB}{245,246,248}

\titleformat{\section}{\large\bfseries\color{mainblue}}{}{0em}{}[\titlerule]
\titleformat{\subsection}{\normalsize\bfseries\color{darkgray}}{}{0em}{}
\titlespacing*{\section}{0pt}{0.9em}{0.45em}
\titlespacing*{\subsection}{0pt}{0.65em}{0.25em}

\setlist[enumerate]{leftmargin=2.2em,itemsep=0.25em,topsep=0.25em}
\setlist[itemize]{leftmargin=2.0em,itemsep=0.2em,topsep=0.2em}

\pagestyle{fancy}
\fancyhf{}
\lhead{\small 格物助航｜每日一练}
\rhead{\small 2026.6.11}
\cfoot{\small 第 \thepage\ 页 / 共 \pageref{LastPage} 页}
\renewcommand{\headrulewidth}{0.4pt}

\newtcolorbox{problemBox}{
  colback=lightblue,
  colframe=mainblue,
  arc=2mm,
  boxrule=0.7pt,
  left=1.0em,
  right=1.0em,
  top=0.8em,
  bottom=0.8em,
  before skip=0.6em,
  after skip=0.8em
}

\newtcolorbox{tipBox}{
  colback=lightgray,
  colframe=darkgray,
  arc=1.5mm,
  boxrule=0.5pt,
  left=0.9em,
  right=0.9em,
  top=0.6em,
  bottom=0.6em,
  before skip=0.5em,
  after skip=0.5em
}

\newcommand{\dd}{\mathrm{d}}
\newcommand{\ii}{\bm{i}}
\newcommand{\jj}{\bm{j}}

\begin{document}

\begin{center}
  {\Large\bfseries\color{mainblue} 格物助航｜每日一练}\par
  \vspace{0.25em}
  {\large 力学第六章：质心力学定理}\par
  \vspace{0.35em}
  {\small 2026年6月11日}\par
  \vspace{0.35em}
  {\small 编写：贡泽轩}
\end{center}

\vspace{-0.3em}

\begin{problemBox}
\noindent\textbf{题目}\quad
质量为 $m$ 的质点 $A$ 与质量为 $M$ 的质点 $B$ 只受彼此之间的万有引力作用。
初始时两者相距很远，可认为 $r=\infty$，且在惯性参考系中 $B$ 静止，$A$ 相对 $B$ 的速度大小为 $v_0$。
若不受万有引力作用，质点 $A$ 将沿其初始速度方向作匀速直线运动；质点 $B$ 到这条直线的垂直距离为 $D$，称为碰撞参数或瞄准距离。
\begin{enumerate}
    \item 求两质点之间的最小距离 $r_{\min}$。
    \item 若 $M\gg m$，可近似认为 $B$ 固定不动。若实验测得最小距离为 $d$，反求 $B$ 的质量 $M$。
    \item 说明为什么在最接近点处不能令质点的速度为零。
\end{enumerate}
\end{problemBox}


\section*{详细解法}

\subsection*{一、把两体问题化为相对运动问题}
两质点只受内力作用，且内力沿两质点连线方向，因此这是一个典型的有心力问题。
设相对位矢为
\[
\bm r=\bm r_A-\bm r_B,
\]
两质点间距离为
\[
r=|\bm r|.
\]
引入约化质量
\[
\mu=\frac{mM}{m+M}.
\]
两体相对运动可等效为一个质量为 $\mu$ 的质点在势能场
\[
U(r)=-\frac{GmM}{r}
\]
中运动。
这里不要分别写 $A$、$B$ 的运动方程，而要改写为相对运动。
这样问题就变成了一个“一个质点在有心引力场中运动”的问题。

\subsection*{二、写出角动量守恒}
万有引力始终沿 $\bm r$ 方向，因此对两体连线方向没有力矩，等效质点的角动量守恒：
\[
L=\mu r^2\dot{\theta}=\mu r v_\perp=\text{常量}.
\]
在 $r=\infty$ 时，相对速度为 $v_0$，碰撞参数为 $D$，所以相对运动的初始角动量为
\[
L=\mu v_0D.
\]
设两质点最接近时的距离为 $r_{\min}$，此时相对速度大小为 $u$。
在最接近点，径向速度为零，即
\[
\dot r=0.
\]
因此此时相对速度完全沿切向，有
\[
L=\mu r_{\min}u.
\]
由角动量守恒得
\[
\mu v_0D=\mu r_{\min}u,
\]
所以
\[
u=\frac{v_0D}{r_{\min}}.
\]

\subsection*{三、写出机械能守恒}
由于万有引力是保守力，机械能守恒。
初始时 $r=\infty$，势能为零，因此相对运动的总机械能为
\[
E=\frac{1}{2}\mu v_0^2.
\]
在最近点处，距离为 $r_{\min}$，速度大小为 $u$，势能为
\[
U(r_{\min})=-\frac{GmM}{r_{\min}}.
\]
于是
\[
E=\frac{1}{2}\mu u^2-\frac{GmM}{r_{\min}}.
\]
代入能量守恒：
\[
\frac{1}{2}\mu v_0^2
=
\frac{1}{2}\mu u^2-\frac{GmM}{r_{\min}}.
\]
再代入
\[
u=\frac{v_0D}{r_{\min}},
\]
得到
\[
\frac{1}{2}\mu v_0^2
=
\frac{1}{2}\mu \frac{v_0^2D^2}{r_{\min}^2}
-\frac{GmM}{r_{\min}}.
\]
两边同乘以 $2/\mu$：
\[
v_0^2
=
\frac{v_0^2D^2}{r_{\min}^2}
-\frac{2GmM}{\mu r_{\min}}.
\]
因为
\[
\frac{mM}{\mu}=m+M,
\]
所以
\[
\frac{GmM}{\mu}=G(m+M).
\]
于是
\[
v_0^2
=
\frac{v_0^2D^2}{r_{\min}^2}
-\frac{2G(m+M)}{r_{\min}}.
\]
整理为关于 $r_{\min}$ 的二次方程：
\[
v_0^2 r_{\min}^2
+2G(m+M)r_{\min}
-v_0^2D^2=0.
\]
取正根，得
\[
\boxed{
r_{\min}
=
\frac{
\sqrt{G^2(m+M)^2+v_0^4D^2}
-G(m+M)
}{v_0^2}
}.
\]

\subsection*{四、固定中心近似下反求 $M$}
若 $M\gg m$，则
\[
m+M\simeq M.
\]
若已知最近距离为 $d$，则由上面的二次方程
\[
v_0^2d^2+2GMd-v_0^2D^2=0.
\]
移项得
\[
2GMd=v_0^2(D^2-d^2).
\]
所以
\[
\boxed{
M=\frac{v_0^2(D^2-d^2)}{2Gd}
}.
\]
这也说明：对于吸引力场，轨道会向中心偏折，所以通常有
\[
d<D.
\]
若计算中得到 $d\geq D$，就要检查题意或符号是否出错。

\subsection*{五、说明最近点处速度不是零}
最近点的物理条件是
\[
\dot r=0,
\]
意思是距离 $r$ 不再继续减小，径向速度为零。
但是速度还有切向分量，因此一般有
\[
u\neq 0.
\]
如果误以为最近点速度为零，就会得到
\[
L=\mu r_{\min}u=0,
\]
这与初始角动量
\[
L=\mu v_0D
\]
矛盾。只要 $D\neq 0$，角动量就不为零，最近点处速度就不可能为零。

\section*{技巧与知识点}

\begin{tipBox}
\begin{enumerate}
    \item \textbf{两体问题先转化为相对运动问题。}
    
    对于两个质点之间的相互作用，不要直接硬写两个运动方程。应引入相对位矢
    \[
    \bm r=\bm r_A-\bm r_B
    \]

    和约化质量
    \[
    \mu=\frac{mM}{m+M}.
    \]

    \item \textbf{有心力问题优先寻找守恒量。}
    
    万有引力沿两质点连线方向，因此力矩为零，角动量守恒：
    \[
    L=\mu r v_\perp=\text{常量}.
    \]

    \item \textbf{保守力问题使用机械能守恒。}
    
    引力势能为
    \[
    U(r)=-\frac{GmM}{r}.
    \]

    注意引力势能是负的，这是本题最容易错的符号之一。

    \item \textbf{最近距离的条件是径向速度为零。}
    
    最近点处不是总速度为零，而是
    \[
    \dot r=0.
    \]

    此时速度完全沿切向。

    \item \textbf{碰撞参数 $D$ 对应初始角动量。}
    
    初始角动量不是 $\mu r v_0$，因为初始时 $r=\infty$。
    正确写法是
    \[
    L=\mu v_0D.
    \]

    这里的 $D$ 是速度方向到力心的垂直距离。
\end{enumerate}
\end{tipBox}

\section*{复习建议}

\begin{enumerate}
    \item 复习第六章时，应把“质心”和“相对运动”联系起来看。
    质心描述整体运动，相对坐标描述内部运动。
    
    \item 遇到两体相互作用题，先问自己三个问题：外力是否为零？内力是否为有心力？内力是否为保守力？
    这三个问题分别对应动量守恒、角动量守恒和机械能守恒。

    \item 对有心运动题，建议熟练掌握以下三行：
    \[
    L=\mu r v_\perp=\mu v_0D,
    \]
    \[
    E=\frac{1}{2}\mu v^2+U(r),
    \]
    \[
    \text{最近点：}\quad \dot r=0,\quad v=v_\perp.
    \]
\end{enumerate}

\end{document}
